\chapter{Creio no Espírito Santo}

O oitavo artigo da Profissão de Fé declara: ``Creio no Espírito Santo''. Este artigo confessa a fé na terceira Pessoa da Santíssima Trindade, consubstancial ao Pai e ao Filho, que procede de ambos. O Catecismo da Igreja Católica, nos parágrafos 683 a 747, apresenta o Espírito Santo como Aquele que completa a obra da salvação, santifica a Igreja e habita nos corações dos fiéis. Ele é o dom do Pai e do Filho, o Paráclito prometido por Jesus, que guia a Igreja na verdade e a renova continuamente. A fé no Espírito Santo é inseparável da fé no Pai e no Filho.

\section{A Pessoa do Espírito Santo (683--701)}

A fé no Espírito Santo é possível apenas pela ação do próprio Espírito. Ninguém pode dizer ``Jesus é Senhor'' senão pelo Espírito Santo. Ele é o primeiro a despertar a fé em nós e a comunicar a vida nova. O Espírito Santo foi revelado progressivamente: o Antigo Testamento proclamou claramente o Pai, o Novo Testamento revelou o Filho e deu um vislumbre da divindade do Espírito. Agora, o Espírito habita entre nós e concede-nos uma visão mais clara de si mesmo.

O Espírito Santo é uma das Pessoas da Santíssima Trindade, consubstancial ao Pai e ao Filho. É ``adorado e glorificado com o Pai e o Filho''. Ele procede do Pai e do Filho como de um só princípio. Não é gerado nem não-gerado, mas procede de ambos. É o Amor e a Santidade do Pai e do Filho. Como declara o Símbolo de Fé: ``Creio também no Espírito Santo, completo e perfeito e verdadeiro Deus, procedente do Pai e do Filho, co-igual, co-essencial, co-onipotente e co-eterno com o Pai e o Filho''.

Diversos títulos revelam a sua identidade: Espírito de verdade, Paráclito (Consolador e Advogado), Espírito de adoção, Espírito de Cristo, Espírito do Senhor. Os símbolos que o representam são ricos: água (que significa o novo nascimento no Batismo), unção (que exprime a sua ação consagradora), fogo (que simboliza a sua energia transformadora), nuvem e luz, selo, mão, dedo de Deus e pomba.

O Espírito Santo foi enviado pelo Pai e pelo Filho. Jesus prometeu o Paráclito aos Apóstolos: ``Eu rogarei ao Pai e Ele vos dará outro Advogado para estar convosco para sempre, o Espírito da verdade''. Na Última Ceia, Jesus anunciou que o Espírito Santo ensinaria tudo, recordaria tudo o que Ele dissera e guiaria os discípulos para a verdade completa. No dia de Pentecostes, o Espírito desceu sobre a Igreja nascente sob a forma de línguas de fogo, enchendo os Apóstolos e fazendo-os falar em diversas línguas.

A ação do Espírito Santo é inseparável da de Cristo. Eles têm uma missão conjunta. O Espírito revela Cristo, glorifica-O e faz com que Ele seja conhecido. ``Ele não falará de si mesmo'', mas dirá tudo o que ouvir e anunciará as coisas futuras. É através do Espírito Santo que Cristo continua a estar presente na sua Igreja.

\section{A ação do Espírito Santo na Igreja (702--747)}

O Espírito Santo é a alma da Igreja. Ele habita na Igreja como num templo e nos corações dos fiéis. Ele santifica a Igreja, unifica-a na comunhão e nos ministérios, e distribui os seus dons e carismas para a edificação do Corpo de Cristo. A Igreja é guiada pelo Espírito Santo no caminho de toda a verdade.

Através dos sacramentos, especialmente do Batismo e da Confirmação, o Espírito Santo comunica a vida divina. No Batismo, Ele faz de nós novas criaturas e membros do Corpo de Cristo. Na Confirmação, Ele fortalece a nossa fé e nos sela com o dom da sua força. Ele é o princípio da vida nova que nos torna filhos adotivos do Pai.

O Espírito Santo é o autor da Sagrada Escritura e ilumina a Igreja para a compreender e interpretar fielmente. Ele assiste o Magistério da Igreja para que esta conserve e expõe fielmente o depósito da fé. Ele desperta nos fiéis o sentido sobrenatural da fé, pelo qual o Povo de Deus adere infalivelmente à verdade revelada.

Na vida cristã, o Espírito Santo produz frutos de santidade: caridade, alegria, paz, paciência, benignidade, bondade, longanimidade, mansidão e fidelidade. Ele concede carismas para o serviço da Igreja e inspira a oração dos fiéis, ensinando-nos a clamar ``Abba, Pai!''. Ele intercede por nós com gemidos inefáveis.

A Igreja vive da contínua efusão do Espírito Santo. No Pentecostes, o Espírito desceu sobre Maria e os Apóstolos reunidos em oração, inaugurando publicamente a missão da Igreja. Desde então, o Espírito Santo renova constantemente a Igreja, mantendo-a jovem e fiel à sua missão. Ele é o princípio da comunhão entre os fiéis e o vínculo de caridade que une a Igreja na terra à Igreja do céu.

A fé no Espírito Santo impulsiona os cristãos a viverem em vigilância e docilidade ao seu agir. Ele é o mestre interior que nos ensina a seguir Cristo, a amar o Pai e a servir os irmãos. A invocação do Espírito Santo deve acompanhar a vida cristã diária, especialmente antes da celebração da Eucaristia e nos momentos importantes de decisão. ``Vinde, Espírito Santo'' é a prece constante da Igreja que peregrina rumo à plenitude do Reino.

A confissão ``Creio no Espírito Santo'' compromete-nos a reconhecer a sua presença e a colaborar com a sua graça. Ele é o dom por excelência que o Pai e o Filho nos concedem. Acolhê-lo significa deixar-se conduzir por Ele para a glória do Pai e a salvação do mundo.
